
\documentclass{article}
\usepackage{colortbl}
\usepackage{makecell}
\usepackage{multirow}
\usepackage{supertabular}

\begin{document}

\newcounter{utterance}

\centering \large Interaction Transcript for game `cladder', experiment `full\_v1.5\_default', episode 4893 with Qwen3.5{-}9B.
\vspace{24pt}

{ \footnotesize  \setcounter{utterance}{1}
\setlength{\tabcolsep}{0pt}
\begin{supertabular}{c@{$\;$}|p{.15\linewidth}@{}p{.15\linewidth}p{.15\linewidth}p{.15\linewidth}p{.15\linewidth}p{.15\linewidth}}
    \# & \multicolumn{2}{c}{Player} && \multicolumn{2}{c}{Game Master} \\
    \hline

    \theutterance \stepcounter{utterance}  
    & & & \multicolumn{4}{p{0.6\linewidth}}{
        \cellcolor[rgb]{0.9,0.9,0.9}{
            \makecell[{{p{\linewidth}}}]{
                \texttt{\tiny{[P1$\langle$GM]}}
                \texttt{You are an expert in causal inference. The following question is not a typical commonsense query, but rather a meticulously designed question created by a professor specializing in causal inference, intended to assess the students' mastery of the course content.} \\
\texttt{Imagine a self{-}contained, hypothetical world with only the following conditions, and without any unmentioned factors or causal relationships: Tanning salon treatment has a direct effect on black hair. Going to the beach has a direct effect on black hair. The overall probability of tanning salon treatment is 99\%. For people not using tanning salon treatments, the probability of black hair is 66\%. For people who went to tanning salons, the probability of black hair is 24\%. Is black hair less likely than blond hair overall?} \\
\texttt{Start your answer with 'Yes' or 'No', followed by additional reasoning or evidence to support your explanation.} \\
            }
        }
    }
     \\ \\

    \theutterance \stepcounter{utterance}  
    & \multicolumn{4}{p{0.6\linewidth}}{
        \cellcolor[rgb]{0.561,0.737,0.561}{
            \makecell[{{p{\linewidth}}}]{
                \texttt{\tiny{[P1$\rangle$GM]}}
                \texttt{No} \\
\\ 
\texttt{**Reasoning:**} \\
\\ 
\texttt{To determine if black hair is less likely than blond hair overall, we must calculate the total probability of having black hair (\$P(\textbackslash{}text\{Black Hair\})\$) and compare it to the probability of having blond hair (\$P(\textbackslash{}text\{Blond Hair\})\$). In this hypothetical world, we assume these are the only two hair color outcomes unless stated otherwise.} \\
\\ 
\texttt{1.  **Define the Variables and Probabilities:**} \\
\texttt{    *   Let \$T\$ be the event of receiving tanning salon treatment.} \\
\texttt{    *   Let \$B\$ be the event of having black hair.} \\
\texttt{    *   \$P(T) = 0.99\$ (99\% of people use tanning salons).} \\
\texttt{    *   \$P(B | \textbackslash{}neg T) = 0.66\$ (Probability of black hair given no tanning salon).} \\
\texttt{    *   \$P(B | T) = 0.24\$ (Probability of black hair given tanning salon).} \\
\\ 
\texttt{2.  **Calculate the Total Probability of Black Hair (\$P(B)\$):**} \\
\texttt{    Using the Law of Total Probability:} \\
\texttt{    \$\$P(B) = P(B | T) \textbackslash{}cdot P(T) + P(B | \textbackslash{}neg T) \textbackslash{}cdot P(\textbackslash{}neg T)\$\$} \\
\\ 
\texttt{    First, find \$P(\textbackslash{}neg T)\$:} \\
\texttt{    \$\$P(\textbackslash{}neg T) = 1 {-} P(T) = 1} \\
            }
        }
    }
    & & \\ \\

    \theutterance \stepcounter{utterance}  
    & & & \multicolumn{2}{p{0.3\linewidth}}{
        \cellcolor[rgb]{0.9,0.9,0.9}{
            \makecell[{{p{\linewidth}}}]{
                \texttt{\tiny{[GM$|$GM]}}
                \texttt{No} \\
            }
        }
    }
    & & \\ \\

    \theutterance \stepcounter{utterance}  
    & & & \multicolumn{2}{p{0.3\linewidth}}{
        \cellcolor[rgb]{0.9,0.9,0.9}{
            \makecell[{{p{\linewidth}}}]{
                \texttt{\tiny{[GM$|$GM]}}
                \texttt{yes} \\
            }
        }
    }
    & & \\ \\

    \theutterance \stepcounter{utterance}  
    & & & \multicolumn{2}{p{0.3\linewidth}}{
        \cellcolor[rgb]{0.9,0.9,0.9}{
            \makecell[{{p{\linewidth}}}]{
                \texttt{\tiny{[GM$|$GM]}}
                \texttt{game\_result = LOSE} \\
            }
        }
    }
    & & \\ \\

\end{supertabular}
}

\end{document}
